% Ad M. Brutum 1.16 (ad-brutum-01-16)
% Source: https://raw.githubusercontent.com/PerseusDL/canonical-latinLit/master/data/phi0474/phi059/phi0474.phi059.perseus-lat1.xml
% Fetched by scripts/fetch_latin.py

% Dateline (Perseus): Scr. Athenis in. Maio a. 711 (a).

\ciceroLetterOpener{Brutus Ciceroni salutem}

\ciceroSection{1} particulam litterularum tuarum, quas misisti Octavio, legi missam ab Attico mihi. studium tuum curaque de salute mea nulla me nova voluptate adfecit. non solum enim usitatum sed etiam cotidianum est aliquid audire de te, quod pro nostra dignitate fideliter atque honorifice dixeris aut feceris. at dolore quantum maximum capere animo possum eadem illa pars epistulae scripta ad Octavium de nobis adfecit. sic enim illi gratias agis de re publica, tam suppliciter ac demisse—quid scribam? pudet condicionis ac fortunae sed tamen scribendum est: commendas nostram salutem illi, quae morte qua non perniciosior? ut prorsus prae te feras non sublatam dominationem sed dominum commutatum esse. verba tua recognosce et aude negare servientis adversus regem istas esse preces. Vnum ais esse quod ab eo postuletur et exspectetur, ut eos civis de quibus viri boni populusque Romanus bene existimet salvos velit. quid si nolit? non erimus? atqui non esse quam esse per illum praestat.

\ciceroSection{2} ego medius fidius non existimo tam omnis deos aversos esse a salute populi Romani ut Octavius orandus sit pro salute cuiusquam civis, non dicam pro liberatoribus orbis terrarum; iuvat enim magnifice loqui et certe decet adversus ignorantis quid pro quoque timendum aut a quoque petendum sit. hoc tu, Cicero, posse fateris Octavium et illi amicus es? aut, si me carum habes, vis Romae videre, cum ut ibi esse possem commendandus puero illi fuerim? cui quid agis gratias, si ut nos salvos esse velit et patiatur rogandum putas? an hoc pro beneficio habendum est, quod se quam Antonium esse maluerit a quo ista petenda essent? Vindici quidem alienae dominationis, non vicario, ecquis supplicat ut optime meritis de re publica liceat esse salvis?

\ciceroSection{3} ista vero imbecillitas et desperatio, cuius culpa non magis in te residet quam in omnibus aliis, et Caesarem in cupiditatem regni impulit et Antonio post interitum illius persuasit ut interfecti locum occupare conaretur et nunc puerum istum extulit, ut tu iudicares precibus esse impetrandam salutem talibus viris misericordiaque unius vix etiam nunc viri tutos fore nos, haud ulla alia re. quod si Romanos nos esse meminissemus, non audacius dominari cuperent postremi homines quam id nos prohiberemus, neque magis inritatus esset Antonius regno Caesaris quam ob eiusdem mortem deterritus.

\ciceroSection{4} tu quidem consularis et tantorum scelerum vindex, quibus oppressis vereor ne in breve tempus dilata sit abs te pernicies, qui potes intueri quae gesseris, simul et ista vel probare vel ita demisse ac facile pati ut probantis speciem habeas? quod autem tibi cum Antonio privatim odium? nempe quia postulabat haec, salutem ab se peti, precariam nos incolumitatem habere a quibus ipse libertatem accepisset, esse arbitrium suum de re publica, quaerenda esse arma putasti quibus dominari prohiberetur, scilicet ut illo prohibito rogaremus alterum qui se in eius locum reponi pateretur, an ut esset sui iuris ac mancipi res publica? nisi forte non de servitute sed de condicione serviendi recusatum est a nobis. atqui non solum bono domino potuimus Antonio tolerare nostram fortunam sed etiam beneficiis atque honoribus ut participes frui quantis vellemus. quid enim negaret iis quorum patientiam videret maximum dominationis suae praesidium esse? sed nihil tanti fuit quo venderemus fidem nostram et libertatem.

\ciceroSection{5} hic ipse puer quem Caesaris nomen incitare videtur in Caesaris interfectores, quanti aestimet, si sit commercio locus, posse nobis auctoribus tantum quantum profecto potent, quoniam vivere et pecunias habere et dici consulares volumus! ceterum ne nequiquam perierit ille cuius interitu quid gavisi sumus, si mortuo nihilo minus servituri eramus, nulla cura adhibetur? sed mihi prius omnia di deaeque eripuerint quam illud iudicium, quo non modo heredi eius quem occidi non concesserim quod in illo non tuli, sed ne patri quidem meo, si revivescat, ut patiente me plus legibus ac senatu possit. an hoc tibi persuasum est, fore ceteros ab eo liberos quo invito nobis in ista civitate locus non sit? qui porro id quod petis fleri potest ut impetres? rogas enim velit nos salvos esse. videmur ergo tibi salutem accepturi cum vitam acceperimus? quam , nisi prius dimittimus dignitatem et libertatem, qui possumus accipere?

\ciceroSection{6} an tu Romae habitare, id putas incolumem esse? res non locus oportet praestet istuc mihi. neque incolumis Caesare vivo fui, nisi postea quam illud conscivi facinus, neque usquam exsul esse possum, dum servire et pati contumelias peius odero malis omnibus aliis. nonne hoc est in easdem tenebras recidisse, si ab eo qui tyranni nomen adscivit sibi, cum in Graecis civitatibus liberi tyrannorum oppressis illis eodem supplicio adficiantur, petitur ut vindices atque oppressores dominationis salvi sint? hanc ego civitatem videre velim aut putem ullam, quae ne traditam quidem atque inculcatam libertatem recipere possit plusque timeat in puero nomen sublati regis quam confidat sibi, cum illum ipsum qui maximas opes habuerit paucorum virtute sublatum videat? me vero posthac ne commendaveris Caesari tuo, ne te quidem ipsum, si me audies. valde care aestimas tot annos quot ista aetas recipit, si propter eam causam puero isti supplicaturus es.

\ciceroSection{7} deinde quod pulcherrime fecisti ac facis in Antonio vide ne convertatur a laude maximi animi ad opinionem formidinis. nam si Octavius tibi placet, a quo de nostra salute petendum sit, non dominum fugisse sed amiciorem dominum quaesisse videberis. quem quod laudas ob ea quae adhuc fecit plane probo; sunt enim laudanda, si modo contra alienam potentiam non pro sua suscepit eas actiones. Cum vero iudicas tantum illi non modo licere sed etiam a te ipso tribuendum esse ut rogandus sit ne nolit esse nos salvos, nimium magnam mercedem statuis (id enim ipsum illi largiris quod per illum habere videbatur res publica), neque hoc tibi in mentem venit, si Octavius ullis dignus sit honoribus quia cum Antonio bellum gerat, iis qui illud malum exciderint cuius istae reliquiae sunt nihil quo expleri possit eorum meritum tributurum umquam populum Romanum, si omnia simul congesserit.

\ciceroSection{8} ac vide quanto diligentius homines metuant quam meminerint, quia Antonius vivat atque in armis sit, de Caesare vero quod fieri potuit ac debuit transactum est neque iam revocari in integrum potest. Octavius is est qui quid de nobis iudicaturus sit exspectet populus Romanus; nos ii sumus de quorum salute unus homo rogandus videatur? ego vero, ut istoc revertar, is sum qui non modo non supplicem sed etiam coerceam postulantis ut sibi supplicetur. aut longe a servientibus abero mihique esse iudicabo Romam ubicumque liberum esse licebit, ac vestri miserebor quibus nec aetas neque honores nec virtus aliena dulcedinem vivendi minuere potuerit.

\ciceroSection{9} mihi quidem ita beatus esse videbor, si modo constanter ac perpetuo placebit hoc consilium ut relatam putem gratiam pietati meae. quid enim est melius quam memoria recte factorum et libertate contentum neglegere humana? sed certe non succumbam succumbentibus nec vincar ab iis qui se vinci volunt experiarque et temptabo omnia neque desistam abstrahere a servitio civitatem nostram. si secuta fuerit quae debet fortuna, gaudebimus omnes; si minus, ego tamen gaudebo. quibus enim potius haec vita factis aut cogitationibus traducatur quam iis quae pertinuerint ad liberandos civis meos?

\ciceroSection{10} te , Cicero, rogo atque hortor ne defetigere neu diffidas, semper in praesentibus malis prohibendis futura quoque explores ne se, nisi ante sit occursum, insinuent. fortem et liberum animum, quo et consul et nunc consularis rem publicam vindicasti, sine constantia et aequabilitate nullum esse putaris. fateor enim duriorem esse condicionem spectatae virtutis quam incognitae. bene facta pro debitis exigimus, quae aliter eveniunt ut decepti ab iis infesto animo reprehendimus. itaque resistere Antonio Ciceronem, etsi maxima laude dignum est, tamen, quia ille consul hunc consularem merito praestare videtur, nemo admiratur;

\ciceroSection{11} idem Cicero, si flexerit adversus alios iudicium suum quod tanta firmitate ac magnitudine direxit in exturbando Antonio, non modo reliqui temporis gloriam eripuerit sibi sed etiam praeterita evanescere coget. nihil enim per se amplum est nisi in quo iudici ratio exstat. quin neminem magis decet rem publicam amare libertatisque defensorem esse vel ingenio vel rebus gestis vel studio atque efflagitatione omnium. qua re non Octavius est rogandus ut velit nos salvos esse, magis tute te exsuscita, ut eam civitatem in qua maxima gessisti liberam atque honestam fore putes, si modo sint populo duces ad resistendum improborum consiliis.
